%% ma: L10-B12-0050
%% lop: 10
%% chuong: 5
%% bai: 12
%% dang: 10.12.8
%% loai: TN4
%% muc: TH
%% dap_an: "C"
%% nguon: "chuyen-de-cac-so-dac-trung-cua-mau-so-lieu-khong-ghep-nhom-toan-10-chuong-trinh-moi.pdf (GV. Phan Nhật Linh) (P01), Bài 01 Phần I Câu 33"
%% kien_thuc: "Lớp 10 Bài 12 (số gần đúng, sai số tuyệt đối, độ chính xác, sai số tương đối, quy tròn số)"
%% trang_thai: chinh-thuc
%% ghi_chu: "Tự duyệt theo yêu cầu giáo viên 22/10/2026."
\begin{ex}
Đường kính của một đồng hồ cát là $8{,}52$ m. Dùng giá trị gần đúng của $\pi=3{,}141592654\ldots$ làm tròn đến hàng phần trăm để tính chu vi của đồng hồ. Kết quả chu vi làm tròn đến hàng phần mười là
\choice{$26{,}75$ m}{$26{,}7$ m}{\True $26{,}8$ m}{$26{,}752$ m}
\loigiai{$\pi\approx3{,}14$ nên $C=\pi d\approx3{,}14\cdot8{,}52=26{,}7528$ (m).

Làm tròn đến hàng phần mười: $26{,}8$ m. Chọn C.}
\end{ex}
